(1) 环上每个电荷元到轴上 z 点的距离均为 \(\sqrt{a^2+z^2}\)，电势线性叠加立即给
\[
\phi(z)=\frac{Q}{4\pi\epsilon_0\sqrt{a^2+z^2}}.
\]
这里无须作角积分，因为距离对环上角度不变。

(2) 轴向电场由导数得
\[
E_z=-\phi'(z)=\frac{Qz}{4\pi\epsilon_0(a^2+z^2)^{3/2}}.
\]
测试粒子受力 \(F_z=qE_z\)，平衡位置 z=0；附近线性力 \(F_z\simeq qQz/(4\pi\epsilon_0a^3)\)，当 qQ<0 时回复稳定，qQ>0 时排斥不稳定。

(3) 对稳定异号电荷 q<0，牛顿方程近似为 \(m\ddot z=-|q|Qz/(4\pi\epsilon_0a^3)\)，所以小振动角频率
\[
\omega=\sqrt{|q|Q/(4\pi\epsilon_0ma^3)}.
\]
若 q>0，频率平方为负而出现指数不稳定，不能把其绝对值称为振动频率。

(4) 对 \(z\gg a\) 且 z>0，使用二项展开
\[
\phi(z)=\frac{Q}{4\pi\epsilon_0z}(1+a^2/z^2)^{-1/2}
=\frac{Q}{4\pi\epsilon_0z}\left[1-\frac{a^2}{2z^2}+O(a^4/z^4)\right].
\]
首项是总净电荷的单极场；因对称性偶极矩为零，下一项为四极阶修正。
